\subsection{Проектирование пользовательского интерфейса}

Интерфейс спроектирован по принципам раздела~1: фильтры отделены от области результатов, обратная связь при смене параметров --- немедленная. Поддерживаются режимы «Стандарт» (светлая/тёмная/авто), альтернативные цветовые схемы («Зелёная», «Розовая») и отдельная тёмная консоль руководителя.

\subsubsection{Структура основного экрана}

Экран дашборда (роль \texttt{analyst}) делится на шапку и рабочую зону. В шапке --- выбор города, источник данных, переключатель темы, имя пользователя и выход. Ниже --- боковая панель фильтров и основная колонка с KPI, графиками, картой, сравнением сегментов и таблицей.

\subsubsection{Панель фильтров}

Боковая панель на десктопе закреплена слева; на узком экране открывается поверх контента. В верхней части --- счётчик «$N$ из $M$ объектов» и кнопка сброса. Ниже --- слайдеры диапазонов цены и площади, чипы комнатности, переключатели исключения этажей. В исходном состоянии после загрузки снимка диапазоны совпадают с baseline summary, счётчик показывает полную выборку.

Конструктор \textbf{дополнительных условий} позволяет задавать произвольные поля CSV --- тип дома, этажность, наличие мебели и т.д. Пример панели после сужения выборки приведён в разделе~4 (рис.~\ref{fig:filters-after}).

\subsubsection{Обзор рынка и визуализации}

Первый блок основной колонки --- \textbf{KPI}: медианы цены, цены за м², площадей, число объектов, текстовые инсайты.

На рис.~\ref{fig:dash-main-light} --- общий вид шапки и KPI на полной московской выборке.

\screenshot{dash-main-light}{dash-main-light-moscow}{Шапка дашборда, KPI и автоинсайты (Москва, 12\,549 объектов, светлая тема)}

Ниже --- \textbf{сетка графиков}: гистограмма, scatter, столбчатые диаграммы; у аналитика доступен конструктор пользовательских виджетов (рис.~\ref{fig:dash-charts-full}).

\screenshot{dash-charts-full}{dash-charts-light-full}{Блок визуализаций: гистограмма цены и scatter на полной выборке}

\subsubsection{Конструктор графиков}

Sidebar \texttt{ChartSidebar} доступен ролям analyst и admin: быстрые пресеты типовых графиков, выбор числового или категориального столбца для нового виджета, drag-and-drop для изменения порядка виджетов на холсте, экспорт PNG/PDF. Результат работы конструктора на отфильтрованной выборке --- в разделе~4 (рис.~\ref{fig:chart-histogram-detail}, \ref{fig:chart-categories-year}).

\subsubsection{Карта объектов}

Карта на базе JavaScript API Яндекс.Карт~\cite{yandex-maps-api}: кластеры объявлений, балун с ценой и ссылкой на карточку, настройка лимита точек (рис.~\ref{fig:dash-map-light}).

\screenshot{dash-map-light}{dash-map-light-moscow}{Карта Москвы с кластерами и балуном (полная выборка)}

\subsubsection{Сравнение сегментов}

Два независимых набора фильтров и KPI рядом --- типичный сценарий «новостройки vs вторичка» (рис.~\ref{fig:dash-segments-full}).

\screenshot{dash-segments-full}{dash-segments-light-full}{Сравнение сегментов «новостройки» и «вторичка» (полная выборка)}

\subsubsection{Таблица объявлений}

Таблица для ролей analyst и admin: сортировка, настраиваемые колонки, геокодирование строк без адреса, экспорт CSV/XLSX (рис.~\ref{fig:dash-table-full}).

\screenshot[0.95]{dash-table-full}{dash-table-light-full}{Таблица объявлений на полной выборке (12,5~тыс.\ строк)}

\subsubsection{Цветовые темы}

Помимо светлой и тёмной базовой темы доступны альтернативные палитры; зелёная схема для датасета Саратова (75 объектов) (рис.~\ref{fig:theme-green-saratov}).

\screenshot{theme-green-saratov}{dash-theme-green-saratov}{Зелёная цветовая схема (Саратов, 75 объектов)}

\subsubsection{Интерфейсы по ролям}

\textbf{Наблюдатель} (\texttt{observer}) видит KPI, графики, карту и упрощённый список карточек без полной таблицы и экспорта (рис.~\ref{fig:observer-saratov}, розовая схема).

\screenshot{observer-saratov}{observer-saratov-light}{Интерфейс observer: карта и карточки объявлений (Саратов)}

\textbf{Администратор} (\texttt{admin}) вместо блока «Обзор рынка» получает панель серверного датасета: upload CSV, ingest, история загрузок (рис.~\ref{fig:admin-server}).

\screenshot{admin-server}{admin-server}{Панель admin: серверный датасет (upload, ingest, история)}

Отдельно доступна панель запуска парсера с индикатором прогресса задачи (рис.~\ref{fig:admin-parser-panel}).

\screenshot{admin-parser-panel}{admin-parser-running}{Панель admin: запуск парсера и индикатор прогресса}

\textbf{Руководитель} (\texttt{manager}) работает в отдельной консоли. Вкладка «Рынки» --- сводка по городам (рис.~\ref{fig:manager-markets}).

\screenshot{manager-markets}{manager-markets-dark}{Консоль manager: вкладка «Рынки» (тёмная тема)}

Вкладка «Пользователи» --- CRUD учётных записей (рис.~\ref{fig:manager-users}).

\screenshot[0.88]{manager-users}{manager-users-dark}{Консоль manager: вкладка «Пользователи» (тёмная тема)}

\subsubsection{Типовые сценарии}

\textbf{Аналитик:} выбор источника данных $\rightarrow$ настройка фильтров $\rightarrow$ просмотр KPI и графиков $\rightarrow$ карта и таблица $\rightarrow$ экспорт. Проверка сценария с применёнными фильтрами --- раздел~4.

\textbf{Администратор:} upload/ingest $\rightarrow$ запуск стадии парсера $\rightarrow$ при капче --- ручное продолжение $\rightarrow$ ingest готового CSV. Результаты проверки --- раздел~4.

\textbf{Наблюдатель} и \textbf{руководитель} --- сценарии просмотра и управления без изменения датасетов; подтверждение на production --- раздел~4.
